\subsection{Methodisches Vorgehen}
\label{subsec:methodisches-vorgehen}

Die Untersuchung folgt einer technischen Fallstudie
\parencite{runeson2009case-study}. Eine Bestandsaufnahme erfasst Struktur,
Komponenten und Verantwortungsgrenzen des Master-Repository. Die anschließende
Analyse bindet den Beobachtungsstand an Branch \texttt{main}, Version
\texttt{0.1.0} und Commit
\path{461fc7519e0db638330904a7496c488e8a0d18bc}; der Arbeitsbaum war zu Beginn
sauber. Der Governance-Commit wird seinem Parent \texttt{0cbe1ba} in einem
Vorher-Nachher-Vergleich gegenübergestellt
\parencite{kleiveist2026governance-snapshot,kleiveist2026governance-commit}.

Die Architekturanalyse betrachtet Frontend, Backend, Desktop-Integration,
Schnittstellen, Persistenz und Deployment. Sie trennt belegte Eigenschaften
von Annahmen und offenen Prüfpunkten \parencite{kleiveist2026architecture}.

Die Profilanalyse untersucht gemeinsamen Kern, Varianten und optionale
Capabilities. Ergänzend wird das Python-Tooling entlang von Generierung,
Systemprüfung, Installation, Entwicklung, Tests, Builds und Release betrachtet
\parencite{kleiveist2026profiles,kleiveist2026tooling}.
Projektinterne Architektur-, Profil- und Werkzeugbeschreibungen dienen dabei
als Primärquellen für den dokumentierten Stand.

Für die Governance werden Konfiguration, Implementierung, Grenzwert- und
Regressionstests, CLI-Exitcodes, Workflow-Dateien, lokale Kommandoausgaben und
aktuelle CI-Läufe abgeglichen. Eine interne Beweismatrix ordnet jeder zentralen
Behauptung Repository-, Test- und gegebenenfalls CI-Belege sowie ein Zielkapitel
zu. Die Aussage, dass 901 Code-LOC regulär einen Fehler auslösen, beruht damit
auf Konfiguration, Klassifikation, dem Test des Übergangs 900/901 und der
nicht ignorierten Rückgabe des Quality-Befehls; festgestellte Gegenbeispiele
werden als Einschränkung ausgewiesen.

Am 22.~August~2026 wurden Hilfebefehle, fokussierte und vollständige
Quality-Läufe, 333 direkt ausführbare Tooling- und Fallstudientests sowie die
öffentlichen Pfade \texttt{doctor}, \texttt{test --suite all},
\texttt{build web} und \texttt{docs index --dry-run} ausgeführt.
Die technische Basis blieb dabei Commit \texttt{461fc751}; die
Fallstudientests prüften zugleich den entstehenden, noch nicht versionierten
deutschen Dokumentstand. Fehlende
Systembibliotheken, Dienste oder Toolchains werden weder als Erfolg noch als
Test des deaktivierten Gegenstands gewertet. Externe oder produktspezifische
Prüfungen bleiben getrennt. Diese Evidenz ergänzt die frühere Abnahme
\parencite{kleiveist2026acceptance}. Abbildung~\ref{fig:case-study-method}
zeigt die Abfolge der Untersuchung.
Bewertung und Transfer folgen damit erst nach Bestandsaufnahme, Analyse und
Verifikation; offene Prüfpunkte werden nicht als bestätigte Eigenschaften
behandelt.

\begin{figure}[H]
  \centering
  \scalebox{1.15}{%
  \begin{tikzpicture}[node distance=6mm]
    \node[casePrimary,minimum width=78mm] (capture) {Snapshot und Governance-Diff erfassen};
    \node[caseBox,minimum width=78mm,below=of capture] (architecture) {Architektur analysieren};
    \node[caseBox,minimum width=78mm,below=of architecture] (profiles) {Profilmodell untersuchen};
    \node[caseBox,minimum width=78mm,below=of profiles] (tooling) {Tooling und Governance untersuchen};
    \node[caseBox,minimum width=78mm,below=of tooling] (verification) {Konfiguration, Tests und CI triangulieren};
    \node[caseOptional,minimum width=78mm,below=of verification] (evaluation) {Stärken und Grenzen bewerten};
    \node[caseOptional,minimum width=78mm,below=of evaluation] (transfer) {Einsatzmöglichkeiten ableiten};

    \draw[caseFlow] (capture) -- (architecture);
    \draw[caseFlow] (architecture) -- (profiles);
    \draw[caseFlow] (profiles) -- (tooling);
    \draw[caseFlow] (tooling) -- (verification);
    \draw[caseFlow] (verification) -- (evaluation);
    \draw[caseFlow] (evaluation) -- (transfer);
  \end{tikzpicture}}
  \caption{Methodische Schritte der Fallstudie}
  \label{fig:case-study-method}
  \smallskip
  {\footnotesize Quelle: Eigene Darstellung.\par}
\end{figure}

\beginnerinput{workspace/statement/100_einleitung/150}
